\begin{textsample}{POS Dim 4 – 2021\_c6 – Score 1.00 – 2021\_c6/enade\_2021\_mat\_6642fe63dd.txt}  \label{ex:f4_pos_069}
CAPÍTULO 6 ANÁLISE TÉCNICA DA PROVA Este capítulo tem por objetivo apresentar o desempenho dos estudantes concluintes de Licenciatura em Matemática no Enade/2021.
Para isso, foram calculadas as estatísticas básicas da prova em seu todo ( seção 6.1.1 ), bem como as estatísticas relacionadas aos Componentes de Formação Geral ( seção 6.1.2 ) e de Conhecimento Específico da Área ( seção 6.1.3 ).
Dadas as suas características, foram analisadas, em separado, as questões objetivas ( seção 6.2 ) e as questões discursivas ( seção 6.3 ).
Tomando -se como base as duas Questões Discursivas do Componente de Formação Geral, nas seções 6.3.1.6 a 6.3.1.8, são apresentados os resultados e os comentários sobre a correção das respostas em relação à Língua Portuguesa.
Para as questões objetivas, foram disponibilizados os Índices de Facilidade e de Discriminação Ponto-Bisserial, também em separado, para os Componentes de Formação Geral ( seção 6.2.1 ) e de Conhecimento Específico ( seção 6.2.2 ).
De cada componente, uma das questões foi escolhida para exemplificar a análise gráfica, relacionando -se as alternativas escolhidas pelos estudantes ( inclusive o gabarito ) com o número de acertos no componente.
No Anexo I, apresenta -se a íntegra da análise gráfica para todas as questões objetivas.
Para cada uma das questões discursivas, os conteúdos dos tipos mais comuns de respostas dos estudantes são apresentados e comparados com o padrão de resposta esperado ( ver Anexo IX com o padrão de respostas ).
Nas tabelas, constam as seguintes estatísticas das notas19 : média do desempenho na prova, erro padrão da média, desvio padrão, nota mínima, mediana e nota máxima para cada um de seus componentes.
Tais estatísticas contemplam o total de estudantes concluintes da Área de Licenciatura em Matemática inscritos e presentes à prova do Enade/2021, tendo em vista agregações, ou por Grandes Regiões e o país como um todo, ou por Categoria Administrativa, Organização Acadêmica da IES e Modalidade de Ensino.
Em relação aos gráficos de distribuição de notas, o intervalo considerado foi de 10 unidades, aberto à esquerda e fechado à direita, com exceção do primeiro intervalo, [ 0 ; 10 ], fechado em ambos os extremos.
Para os gráficos de distribuição das notas das questões discursivas, foram consideradas mais duas categorias : questão em branco20 e nota zero. 19 Uma definição dessas estatísticas pode ser encontrada no Glossário. 20 Nesse grupo estão incluídas também as respostas classificadas como nulas ou desconsideradas. 6.1 ESTATÍSTICAS BÁSICAS DA PROVA Esta seção apresenta estatísticas selecionadas e histogramas da nota geral ( 6.1.1 ) e de cada componente : Formação Geral ( 6.1.2 ) e Conhecimento Específico ( 6.1.3 ).
São também apresentadas estatísticas selecionadas de subpopulações, caracterizadas por Grande Região, Categoria Administrativa, Organização Acadêmica e Modalidade de Ensino. 6.1.1 ESTATÍSTICAS BÁSICAS GERAIS Na Tabela 6.1, são apresentadas as Estatísticas Básicas da prova, por Grande Região, dos estudantes concluintes de Licenciatura em Matemática.
A Média das notas da prova, como um todo ( nas seções seguintes serão analisados os Componentes de Formação Geral e de Conhecimento Específico ), foi 39,3, sendo que os estudantes da região Norte obtiveram a Média mais baixa ( 33,9 ), e os da região Sudeste obtiveram a Média mais alta ( 41,6 ).
As demais Médias foram : 38,4, na região Nordeste, 38,9, na região Sul e 38,0, na região Centro-Oeste.
O Desvio padrão para o Brasil, como um todo, foi 14,7, sendo o maior Desvio padrão encontrado na região Sudeste ( 15,4 ), e o menor, na região Norte ( 11,7 ), indicando uma dispersão um pouco menor das notas dessa última região.
A região que obteve a maior nota Máxima foi a Sudeste ( 90,4 ), ao passo que a região que atingiu a menor nota Máxima foi a Norte ( 71,0 ).
A Mediana do Brasil, como um todo, foi 38,1, sendo a maior Mediana obtida na região Sudeste ( 40,9 ), e a menor obtida na região Norte ( 33,5 ).
A nota Mínima do Brasil foi zero, obtida em todas as regiões.
Considerando -se as notas segundo Grande Região, observa -se que existe diferença estatisticamente significativa ao nível de 95 \% entre as Médias das regiões Norte e Sudeste que obtiveram a menor e a maior Média, respectivamente.
Entre as demais regiões as diferenças entre as Médias não são estatisticamente significativas.
O comportamento das notas dos estudantes de todo o Brasil pode ser observado no Gráfico 6.1, que apresenta um histograma com a distribuição das mesmas : uma distribuição unimodal com a moda no intervalo ( 30 ; 40 ].
Na Tabela 6.2, são apresentadas informações referentes à Média da nota final, desagregadas por Categoria Administrativa, por Organização Acadêmica e por Modalidade de Ensino.
A Média das notas da prova, como um todo, foi 39,3.
Em relação a Categoria Administrativa, os estudantes das IES Públicas obtiveram Média mais alta ( 41,1 ) que a Média nacional, e os das IES Privadas obtiveram Média mais baixa ( 36,2 ).
Observa -se que existe diferença estatisticamente significativa entre as Médias das notas das IES Públicas e as das Privadas.
A diferença entre as Médias das regiões Sudeste e Norte ( 7,7 ), a maior e a menor Média, é maior do que a diferença entre a das IES Públicas e Privadas ( 4,9 ), caracterizando uma maior diversidade regional do que administrativa.
O Desvio padrão para as IES Públicas ( 14,9 ) foi ligeiramente superior ao do Brasil, como um todo ( 14,7 ), não caracterizando diferença na dispersão das notas nesta Categoria Administrativa em relação ao país como um todo.
No tocante à Organização Acadêmica, as Universidades ( 40,0 ) e os CEFET/IF ( 40,5 ) obtiveram Média mais alta que a nacional ( 39,3 ).
A Média dos Centros Universitários ( 36,8 ) e a das Faculdades ( 33,4 ) foram menores do que a nacional.
Constata -se que não existe diferença estatisticamente significativa, ao nível de 95 \%, entre as Médias mais altas, de Universidades e CEFET/IF.
Já a diferença entre as Médias mais altas e as das outras duas Organizações Acadêmicas ( Centros Universitários e Faculdades ) têm diferenças estatisticamente significativas.
A Média da Modalidade Educação a Distância ( 35,8 ) foi inferior à Média da Educação Presencial ( 42,5 ), e a diferença entre elas é estatisticamente significativa, ao nível de 95 \%. 6.1.2 ESTATÍSTICAS BÁSICAS NO COMPONENTE DE FORMAÇÃO GERAL Na Tabela 6.3, são apresentadas as Estatísticas Básicas em relação ao componente da prova que avalia a Formação Geral dos estudantes concluintes.
A nota de Formação Geral é obtida levando -se em conta os acertos nas oito questões objetivas e a nota média obtida nas Questões Discursivas 1 e 2, sendo a nota de cada questão discursiva obtida levando -se em conta 80 \% da nota relativa ao conteúdo e 20 \% da nota relativa ao desempenho linguístico.
Os estudantes de todo o Brasil obtiveram desempenho médio de 34,0.
Quanto à variabilidade, o Desvio padrão das notas dos estudantes do Brasil foi 14,9.
A maior Média foi obtida na região Sudeste ( 35,8 ) e a menor, na região Norte ( 30,5 ).
As demais Médias foram : 32,8, na região Nordeste, 34,2, na região Sul e 32,9, na região Centro-Oeste.
Da mesma forma, o maior Desvio padrão foi obtido na região Sudeste ( 15,3 ), e o menor, na região Norte ( 13,6 ).
Os demais desvios foram : 14,6, na região Nordeste, 14,9, na região Sul e 14,5, na região Centro-Oeste.
A maior nota no Componente de Formação Geral da prova do Enade/2021 foi 89,2, obtida por, pelo menos, um estudante na região Sudeste.
A menor nota Máxima foi obtida na região Norte ( 74,3 ).
A Mediana do Brasil, como um todo, foi 33,5, sendo a menor Mediana encontrada na região Norte ( 30,0 ), e a maior, na região Sudeste ( 35,8 ).
A nota Mínima nessa parte foi zero em todas as regiões.
Considerando -se as notas, segundo Grande Região, observa -se que não existe diferença estatisticamente significativa entre as Médias das notas do Componente de Formação Geral obtidas nas regiões Nordeste e Centro-Oeste.
As diferenças entre as demais Médias são estatisticamente significativas.
No Gráfico 6.2, é possível verificar a avaliação do desempenho dos estudantes do Componente de Formação Geral, a partir do histograma da distribuição das notas correspondentes.
A distribuição é unimodal, com moda em ( 20 ; 30 ], intervalo imediatamente anterior ao modal da distribuição de notas da prova, como um todo ( Gráfico 6.1 ).
Nota -se, ainda, que há pouca diferença na dispersão dos dados comparando -se o Gráfico 6.2 com o Gráfico 6.1 ( distribuição das notas da prova ), o que se confirma pela comparação dos Desvios Padrão : 14,7 para a nota da prova, como um todo, e 14,9 para o Componente de Formação Geral.
Na Tabela 6.4, são apresentadas as informações referentes ao desempenho dos concluintes do Componente de Formação Geral, em diferentes agregações : Categoria Administrativa, Organização Acadêmica e Modalidade de Ensino.
Observa -se que existe diferença estatisticamente significativa entre as Médias dos tipos de Categoria Administrativa : a maior Média obtida por estudantes de IES Públicas ( 34,9 ) e a menor, pelos de IES Privadas ( 32,5 ).
Nota -se que, também, não há diferença estatisticamente significativa entre as Médias das Universidades ( 34,4 ) e dos CEFET/IF ( 34,3 ).
Já a diferença entre essas Médias e as das demais Organizações Acadêmicas são estatisticamente significativas : 33,0 para os Centros Universitários e 31,1 para as Faculdades, a menor Média.
No que se refere a Modalidade de Ensino, a Educação Presencial apresentou maior Média do que a Educação a Distância ( Médias 35,6 e 32,3, respectivamente ).
Tais Médias apresentam diferença estatisticamente significativa. 6.1.3 ESTATÍSTICAS BÁSICAS DO COMPONENTE DE CONHECIMENTO ESPECÍFICO Na Tabela 6.5, são apresentadas as Estatísticas Básicas referentes ao Componente de Conhecimento Específico da Área de Licenciatura em Matemática.
A nota do componente de Conhecimentos Específicos leva em conta as notas da parte objetiva e a média das notas obtidas nas três questões discursivas da Área.
A Média do desempenho dos estudantes do Brasil, como um todo, foi 41,1.
A maior Média foi obtida na região Sudeste ( 43,5 ), e a menor, na região Norte ( 35,0 ).
As demais Médias foram : 40,2 na região Nordeste, 40,5 na região Sul e 39,7 na região Centro-Oeste.
Quanto à variabilidade das notas, o Desvio padrão do Brasil, como um todo, foi 16,6, sendo o maior Desvio padrão observado na região Sudeste ( 17,4 ) e o menor, na região Norte ( 13,5 ).
Os demais desvios foram : 15,9 na região Nordeste e 16,4 na região Sul e 16,0 na região Centro-Oeste.
A Mediana das notas dos estudantes de todo o Brasil foi 39,5.
A maior Mediana ocorreu na região Sudeste ( 42,7 ), e a menor, na região Norte ( 33,7 ).
As demais Medianas foram : 38,9 na região Nordeste, 39,1 na região Sul e 38,0 na região Centro-Oeste.
A nota Máxima do Brasil, como um todo, foi 96,4, sendo obtida por pelo menos um estudante na região Sudeste.
As demais notas Máximas foram : 78,1 na região Norte, 95,2 na região Nordeste, 92,9 na região Sul e 86,5 na região Centro-Oeste.
A nota Mínima foi zero em toas as regiões.
Observa -se que existe diferença estatisticamente significativa entre as regiões Norte e Sudeste na análise das Médias das notas do Componente de Conhecimento Específico.
As diferenças entre as Médias das demais regiões não são significativas, porém, se comparadas à maior e à menor Médias, as diferenças são significativas.
Assim como nos Gráficos 6.1 e 6.2, no Gráfico 6.3, apresentado a seguir, é possível ser feita uma avaliação do desempenho de concluintes de Licenciatura em Matemática, em relação ao Componente de Conhecimento Específico, com um histograma da distribuição das notas correspondentes.
Esta também é uma distribuição unimodal, e o grupo modal é o ( 30 ; 40 ], o mesmo da prova, como um todo, e acima do grupo modal para a Formação Geral.
Na Tabela 6.6, é apresentada uma comparação dos resultados em relação a Categoria Administrativa, a Organização Acadêmica e a Modalidade de Ensino, agora levando -se em conta o desempenho de estudantes do Componente de Conhecimento Específico da prova.
No que se refere a Organização Acadêmica, a maior Média foi a dos CEFET/IF ( 42,5 ), vindo a seguir a das Universidades ( 41,8 ), a dos Centros Universitários ( 38,0 ) e, por fim, a das Faculdades ( 34,2 ).
O maior Desvio padrão, também maior do que o valor para o Brasil, como um todo, foi o das Universidades ( 16,9 ).
As Universidades também obtiveram a maior nota Máxima ( 96,4 ).
Os CEFET/IF obtiveram nota Máxima 90,9, nos Centros Universitários e nas Faculdades a nota Máxima foi 90,2.
As Medianas foram : 40,2 nas Universidades, 36,7 nos Centros Universitários, 32,8 nas Faculdades e 42,2 nos CEFET/IF, a maior.
A nota Mínima foi zero para todas as Organizações Acadêmicas.
Observa -se que existe diferença estatisticamente significativa, ao nível de 95 \%, entre as Médias das notas do Componente de Conhecimento Específico das Universidades e CEFET/IF e as demais Organizações Acadêmicas.
Já entre estas – Universidades e CEFET/IF – a diferença das Médias não é estatisticamente significativa.
Quanto a Categoria Administrativa, observa -se um comportamento semelhante àquele da prova, como um todo, e da parte de Formação Geral, ou seja, existe diferença estatisticamente significativa entre as Médias das IES Públicas ( 43,2 ) e as das IES Privadas ( 37,4 ).
Neste caso, também, a maior Média foi obtida por estudantes de IES Públicas de ensino.
Quanto ao comportamento das Médias, no que toca a Modalidade de Ensino, há diferença estatisticamente significativa entre a Média da Educação Presencial ( 44,8 ) e a da Educação a Distância ( 37,0 ). 6.2 ANÁLISE DAS QUESTÕES OBJETIVAS Esta seção apresenta estatísticas selecionadas e histogramas das Questões Objetivas de Formação Geral ( 6.2.1 ) e de Conhecimento Específico ( 6.2.2 ).
São, também, apresentadas e comparadas as médias das subpopulações caracterizadas por Grande Região.21 6.2.1 COMPONENTE DE FORMAÇÃO GERAL Na Tabela 6.7, são apresentadas as Estatísticas Básicas relativas às oito Questões Objetivas do Componente da prova que abrange a Formação Geral dos estudantes.
A Média do Brasil foi 39,7.
A menor Média foi encontrada na região Norte ( 36,0 ), e a maior, na região Sudeste ( 41,9 ).
As demais Médias foram : 38,1 na região Nordeste, 39,4 na região Sul e 39,2 na região Centro-Oeste.
O Desvio padrão do Brasil foi 19,8, sendo o maior Desvio padrão encontrado na região Sudeste ( 20,2 ), e o menor Desvio padrão encontrado na região Norte ( 18,9 ).
As Medianas do Brasil, como um todo, e das cinco regiões foram iguais a 37,5.
A nota Máxima 100,0 foi alcançada nas cinco Grandes Regiões.
As notas Mínimas ( 0,0 ) foram iguais para todas as regiões.
Na Tabela 6.8, são apresentados o Índice de Facilidade e o Índice de Discriminação ( Ponto-Bisserial ) para cada uma das Questões Objetivas do Componente de Formação Geral.
Quanto ao Índice de Facilidade, foram usadas as seguintes cores para diferenciar o nível de dificuldade da questão :  Azul para as questões classificadas com índice Muito fácil ( \&gt ; =0,86 ), verde para as questões classificadas com índice Fácil ( 0,61 a 0,85 ), amarelo para as questões classificadas com Médio ( 0,41 a 0,60 ), vermelho para as questões classificadas com Difícil ( 0,16 a 0,40 ) e roxo para as questões classificadas com Muito difícil ( \&lt ; =0,15 ).
Já quanto ao Índice de Discriminação, foram usadas as seguintes cores para qualificar a questão :  As questões classificadas com índice Fraco receberam a cor vermelha ( \&lt ; =0,19 ), as classificadas com Médio receberam a cor amarela ( 0,20 a 0,29 ), as classificadas com Bom receberam a cor verde ( 0,30 a 0,39 ) e as classificadas com Muito bom ( \&gt ; =0,40 ) receberam a cor azul.
As Questões Objetivas do Componente de Formação Geral, segundo o Índice de Facilidade, foram assim avaliadas : nenhuma teve o Índice de Facilidade classificado como Muito fácil ou Fácil por ter índice de acertos superior a 0,60.
Quatro questões foram consideradas com índice de dificuldade Médio, situando -se no intervalo entre 0,41 e 0,60 do Índice de Facilidade, ou seja, houve entre 41,0 \% e 60,0 \% de acertos.
As outras quatro questões apresentaram acertos na faixa entre 16 \% e 40 \%, razão pela qual foram classificadas encontrado na região Sudeste ( 20,2 ), e o menor Desvio padrão encontrado na região Norte ( 18,9 ).
As Medianas do Brasil, como um todo, e das cinco regiões foram iguais a 37,5.
A nota Máxima 100,0 foi alcançada nas cinco Grandes Regiões.
As notas Mínimas ( 0,0 ) foram iguais para todas as regiões.
Na Tabela 6.8, são apresentados o Índice de Facilidade e o Índice de Discriminação ( Ponto-Bisserial ) para cada uma das Questões Objetivas do Componente de Formação Geral.
Quanto ao Índice de Facilidade, foram usadas as seguintes cores para diferenciar o nível de dificuldade da questão :  Azul para as questões classificadas com índice Muito fácil ( \&gt ; =0,86 ), verde para as questões classificadas com índice Fácil ( 0,61 a 0,85 ), amarelo para as questões classificadas com Médio ( 0,41 a 0,60 ), vermelho para as questões classificadas com Difícil ( 0,16 a 0,40 ) e roxo para as questões classificadas com Muito difícil ( \&lt ; =0,15 ).
Já quanto ao Índice de Discriminação, foram usadas as seguintes cores para qualificar a questão :  As questões classificadas com índice Fraco receberam a cor vermelha ( \&lt ; =0,19 ), as classificadas com Médio receberam a cor amarela ( 0,20 a 0,29 ), as classificadas com Bom receberam a cor verde ( 0,30 a 0,39 ) e as classificadas com Muito bom ( \&gt ; =0,40 ) receberam a cor azul.
As Questões Objetivas do Componente de Formação Geral, segundo o Índice de Facilidade, foram assim avaliadas : nenhuma teve o Índice de Facilidade classificado como Muito fácil ou Fácil por ter índice de acertos superior a 0,60.
Quatro questões foram consideradas com índice de dificuldade Médio, situando -se no intervalo entre 0,41 e 0,60 do Índice de Facilidade, ou seja, houve entre 41,0 \% e 60,0 \% de acertos.
As outras quatro questões apresentaram acertos na faixa entre 16 \% e 40 \%, razão pela qual foram classificadas como Difícil.
Nenhuma questão foi classificada na categoria Muito difícil, apresentando menos de 15 \% de acertos.
O Índice de Facilidade variou de 0,19 a 0,58, e o de Discriminação, de 0,23 a 0,49.
Seis questões, quanto ao Índice de Discriminação, foram tidas como Muito bom, uma questão foi tida como Bom e outra como Médio.
Na Tabela 6.9, é apresentada a distribuição das questões, levando -se em conta, simultaneamente, a classificação dos dois índices.
As seis questões que tiveram Índice de Discriminação Muito bom figuraram em dois níveis de dificuldade Médio e Difícil : quatro classificadas na categoria Médio ( as Questões 3, 5, 6 e 7 ) do Índice de Facilidade, e duas, na categoria Difícil ( as Questões 1 e 4 ).
Em particular, a Questão 3 foi a que apresentou o maior poder discriminatório, com índice 0,49, e foi considerada Médio em termos de Facilidade, com uma proporção de 0,51 acertos.
O máximo de acertos foi alcançado pela Questão 7, com um Índice de Facilidade de 0,58.
A Questão 8 foi considerada a mais Difícil, com Índice de Facilidade 0,19.
O Índice de Discriminação foi Médio para a Questão 8 ( 0,32 ).
No Gráfico 6.4, para exemplificar, analisa -se o comportamento da Questão 3 de Formação Geral.
Trata -se de uma questão considerada com nível Médio, em relação à Facilidade, e a que obteve o maior Índice de Discriminação dessa parte da prova, 0,49.
Neste gráfico, cada uma das cinco curvas representa o percentual de respostas em determinada alternativa da questão, em função do número de acertos dos estudantes nessa parte da prova ( Formação Geral/Múltipla Escolha ), antes de possíveis eliminações pelo critério Ponto-Bisserial.
A curva em verde corresponde à alternativa B, a alternativa correta para essa questão.
Observa -se que entre os estudantes com menor número de acertos, nessa parte do exame, a situação mais frequente foi a escolha da alternativa E ( em vermelho ), incorreta.
Por exemplo, entre os estudantes que acertaram duas respostas, 42,2 \% escolheram a alternativa E ( em vermelho ), 31,4 \% escolheram a alternativa B ( em verde, gabarito ), 11,5 \% escolheram a alternativa D ( em laranja ), 8,8 \% escolheram a alternativa A ( em azul ) e 5,7 \% escolheram a alternativa C ( em preto ).
Dentre os que acertaram duas respostas entre as questões de múltipla escolha de Formação Geral, 0,2 \% deixou a questão em branco e 0,2 \% marcou mais de uma alternativa, invalidando a questão.
À medida que o número de acertos aumenta, indicando desempenho melhor nessa parte da prova, aumenta, concomitantemente, a proporção de estudantes que selecionaram a alternativa correta B, atingindo 100 \% para os estudantes com oito acertos.
Essa análise permite verificar como a questão discriminou os grupos de desempenho, justificando -se o alto índice ( 0,49 ) obtido na questão.
Cumpre notar que não é possível inferir deste gráfico nem o Índice de Facilidade, que seria uma média da proporção ponderada pela quantidade de estudantes com cada uma das notas, nem o Índice de Discriminação Ponto-Bisserial, por razão equivalente.
No caso extremo, no qual a grande concentração dos acertos dos estudantes fosse abaixo de quatro, o Índice de Facilidade seria obrigatoriamente abaixo de 70,0 \% ( neste exemplo ).
Caso a concentração fosse em seis acertos ou mais, o índice seria obrigatoriamente acima de 90,0 \%.
Os gráficos relativos às demais questões de Formação Geral constam do Anexo I. 6.2.2 COMPONENTE DE CONHECIMENTO ESPECÍFICO Na Tabela 6.10, são apresentadas as Estatísticas Básicas em relação às Questões Objetivas do Componente de Conhecimento Específico da prova de Licenciatura em Matemática por Grande Região.
Nessa parte da prova, como será discutido ainda nesta seção, três questões não foram usadas no cômputo das notas por terem sido descartadas pelo critério do Ponto-Bisserial.
Assim, as notas foram calculadas com base em 24 das 27 Questões Objetivas válidas de Conhecimentos Específicos.
A Média do Brasil deste componente foi 43,4.
A menor Média foi observada na região Norte ( 36,8 ), e a maior, na região Sudeste ( 46,1 ).
O Desvio padrão de todo o Brasil foi 17,7, sendo o menor Desvio padrão encontrado na região Norte ( 14,6 ), e o maior, na região Sudeste ( 18,6 ).
A Mediana de todo o Brasil foi 41,7, o mesmo valor da Mediana encontrada nas regiões Nordeste e Sul.
Nas regiões Norte ( 35,4 ) e Centro-Oeste ( 37,5 ), o valor foi menor, e, na região Sudeste ( 45,8 ), o valor foi maior.
A nota Máxima da prova foi 100,00, obtida nas Questões Objetivas do Componente de Conhecimento Específico, por, pelo menos, um estudante da região Sudeste.
Em duas regiões o valor foi 95,8, sendo elas : Nordeste e Centro-Oeste.
E as menores notas Máxima ocorreram na região Centro-Oeste ( 91,7 ) e Norte ( 83,3 ).
A nota Mínima nas Questões Objetivas do Componente de Conhecimento Específico foi zero, obtida por pelo menos um estudante de todas as regiões.
Na Tabela 6.11, são apresentados os Índices de Facilidade e de Discriminação ( Ponto-Bisserial ) das Questões Objetivas do Componente de Conhecimento Específico, para os estudantes de Licenciatura em Matemática.
Para facilitar a diferenciação das questões, serão usadas as mesmas cores da Tabela 6.8 para as diferentes classificações dos Índices de Facilidade e de Discriminação.
A partir do Índice de Facilidade obtido, pode -se concluir que a maior parte das questões objetivas da prova foram consideradas Difícil ou Médio : das 27 questões válidas, 12 foram classificadas em cada um destes níveis.
Duas questões foram classificadas como Fácil e uma considerada Muito difícil.
Nenhuma questão foi classificada como Muito fácil.
Já quanto aos Índices de Discriminação das Questões Objetivas do Componente de Conhecimento Específico da prova, obtém -se como resultado a seguinte classificação : 10 das 27 questões foram consideradas boas, enquanto nove delas tiveram Índice de Discriminação Muito bom.
Assim, para 19 das 27 questões, os Índices de Discriminação foram Bom ou Muito bom.
Dentre as demais, cinco delas foram classificadas como Médio, e outras três, como Fraco, sendo 8, por conseguinte, a quantidade de questões nos dois patamares mais baixos de Discriminação.
Constata -se, assim, que a prova – no que se refere ao Componente de Conhecimento Específico – possuía boa capacidade de discriminar entre aqueles que dominam ou não o conteúdo.
O Índice de Facilidade variou de 0,14 a 0,68 e o de Discriminação, de -0,04 a 0,54.
Na Tabela 6.12, é apresentada a distribuição das questões levando -se em conta, simultaneamente, a classificação dos dois índices.
Dentre as questões que alcançaram os maiores Índices de Discriminação, nove delas foram classificadas com Índice de Discriminação Muito bom, as de número 11, 17, 18, 21, 22, 23, 25, 34 e 35.
A de índice mais alto, igual a 0,54, foi a Questão 25 que, quanto ao Índice de Facilidade, foi classificada como Médio, com 57 \% dos estudantes marcando a opção correta.
A Questão de número 30 foi a mais Difícil entre as 27 questões específicas, com baixo Índice de Facilidade, apenas 14 \% de acertos.
Essa questão apresentou poder discriminatório muito baixo, 0,04, o que comprova ter sido a mais Difícil para os estudantes.
Destaca -se, também, a Questão 9, com Índice de Facilidade 0,18, o que, em termos percentuais, corresponde a 18 \% de estudantes tendo respondido acertadamente.
Seu Índice de Discriminação também foi baixo, 0,03.
Tais questões foram, portanto, pelo critério Ponto-Bisserial, consideradas inadequadas.
Por isso, as Questões 9 e 30, além da Questão 13 ( que apresentou Índice de Discriminação negativo ), foram eliminadas do cômputo da nota final.
A título de exemplo das análises do comportamento das questões objetivas, no Gráfico 6.5, analisa -se a Questão 25 do Componente de Conhecimento Específico.
Essa questão foi considerada, pelas respostas dos estudantes avaliados na prova, como uma questão com nível de Facilidade Médio, 0,57, ou seja, 57 \% dos estudantes assinalaram, acertadamente, a opção E correspondente ao gabarito.
Como já comentado, seu Índice de Discriminação foi igual a 0,54, classificado como Muito bom, o maior valor de Discriminação.
Neste gráfico, cada uma das cinco curvas representa o percentual de respostas em determinada alternativa da Questão 25, em função do número de acertos dos estudantes nessa parte da prova, antes de possíveis eliminações de questões pelo critério Ponto-Bisserial.
A alternativa correta E, representada no gráfico pela curva em vermelho, foi escolhida em maiores proporções pelos estudantes com desempenho melhor nessa parte da prova.
Já as alternativas incorretas, também denominadas distratores, foram selecionadas, principalmente, por aqueles com notas mais baixas.
No caso da Licenciatura em Matemática, como se observa no eixo horizontal do Gráfico 6.5, nenhum estudante acertou mais do que 26 questões dentre as 27 questões objetivas, considerando -se o gabarito original.
A eliminação das três questões com Índice de Discriminação Fraco do cômputo da nota é que possibilitou a existência de nota Máxima igual a 100,0 por pelo menos um estudante da região Sudeste ( visto na Tabela 6.10 ).
Observa -se que a soma não é 100 \%, por causa das questões não respondidas ou com mais de uma opção marcada.
Aqueles com nota zero, na sua quase totalidade, marcaram a opção A ( linha azul ).
A proporção de estudantes que selecionaram a resposta correta E aumenta gradativamente, chegando a atingir 100 \% para 24 acertos ou mais, enquanto a proporção dos que escolheram alternativas incorretas decai, a partir de oito acertos, em função do número de acertos nessa parte da prova.

% matched lemmas: 
\end{textsample}
